\MulaiModul{04}{Pemilihan}

\subsection{Tujuan}

\begin{tujuan}
Setelah menyelesaikan modul ini, mahasiswa mampu:
\begin{enumerate}
    \item menulis label dan memakainya sebagai tujuan lompatan;
    \item menjelaskan kerja \texttt{CMP} dan \emph{flag} \texttt{ZF},
          \texttt{CF}, serta \texttt{SF};
    \item membedakan lompatan untuk bilangan tidak bertanda
          (\texttt{JA}, \texttt{JB}) dan bilangan bertanda
          (\texttt{JG}, \texttt{JL});
    \item menyusun percabangan \texttt{if} dan \texttt{if--else}
          dengan \texttt{JMP};
    \item merakit, menjalankan, dan menelusuri percabangan dengan
          emu8086 atau DEBUG.
\end{enumerate}
\end{tujuan}

\subsection{Dasar singkat}

\subsubsection*{Label}

Label adalah nama untuk suatu lokasi di dalam segmen kode. Label
didefinisikan dengan titik dua, misalnya \texttt{LEWATI:}, tetapi
ditulis tanpa titik dua saat menjadi tujuan lompatan. Nama label
harus unik, tidak diawali angka, tidak mengandung spasi, dan panjangnya
paling banyak 31 karakter. Label diperlukan karena lompatan selalu
menuju lokasi tertentu, bukan menuju nomor baris.

\subsubsection*{CMP dan flag}

\texttt{CMP} mengurangkan operand kedua dari operand pertama, tetapi
hasil pengurangan itu dibuang. Kedua operand tidak berubah; yang
berubah hanya \emph{flag}. Karena itu, \texttt{CMP} dipakai untuk
menyiapkan keputusan, bukan untuk menyimpan hasil.

Tiga flag yang dipakai pada modul ini dijelaskan pada
Tabel~\ref{tab:m04-flag}. \texttt{ZF} menandai hasil nol,
\texttt{CF} menandai pinjaman pada pengurangan bilangan tidak
bertanda, dan \texttt{SF} menyalin bit tanda hasil pengurangan.

\begin{table}[H]
    \centering
    \caption{Flag yang dipakai pada Modul 04.}
    \label{tab:m04-flag}
    \small
    \begin{tabular}{@{}lll@{}}
        \toprule
        \textbf{Flag} & \textbf{Bernilai 1 bila} & \textbf{Petunjuk} \\
        \midrule
        \texttt{ZF} & hasil operasi nol & kedua nilai sama \\
        \texttt{CF} & terjadi pinjaman pada pengurangan &
        tujuan lebih kecil, tidak bertanda \\
        \texttt{SF} & bit tanda hasil operasi bernilai 1 &
        hasil bertanda negatif \\
        \bottomrule
    \end{tabular}
\end{table}

\subsubsection*{JMP dan lompatan bersyarat}

\texttt{JMP} memindahkan kendali ke label tujuan tanpa memeriksa
keadaan apa pun. Lompatan bersyarat memeriksa flag: bila syarat
terpenuhi, eksekusi pindah ke label; bila tidak, eksekusi berlanjut
ke instruksi berikutnya. Lompatan yang dipakai pada modul ini
dirangkum pada Tabel~\ref{tab:m04-jcc}.

\begin{table}[H]
    \centering
    \caption{Lompatan bersyarat dan jenis bilangan yang dibandingkan.}
    \label{tab:m04-jcc}
    \small
    \begin{tabular}{@{}lll@{}}
        \toprule
        \textbf{Instruksi} & \textbf{Jenis} & \textbf{Lompat bila} \\
        \midrule
        \texttt{JE}, \texttt{JZ} & kesamaan & tujuan sama dengan sumber \\
        \texttt{JNE}, \texttt{JNZ} & kesamaan & tujuan tidak sama dengan sumber \\
        \texttt{JA} & tidak bertanda & tujuan lebih besar \\
        \texttt{JAE} & tidak bertanda & tujuan lebih besar atau sama \\
        \texttt{JB} & tidak bertanda & tujuan lebih kecil \\
        \texttt{JBE} & tidak bertanda & tujuan lebih kecil atau sama \\
        \texttt{JG} & bertanda & tujuan lebih besar \\
        \texttt{JGE} & bertanda & tujuan lebih besar atau sama \\
        \texttt{JL} & bertanda & tujuan lebih kecil \\
        \texttt{JLE} & bertanda & tujuan lebih kecil atau sama \\
        \bottomrule
    \end{tabular}
\end{table}

Huruf A dan B dipakai untuk bilangan tidak bertanda, sedangkan huruf
G dan L dipakai untuk bilangan bertanda. Lompatan tidak bertanda
membaca \texttt{CF} dan \texttt{ZF}; lompatan bertanda membaca
\texttt{SF} dan \texttt{ZF}. Pola bit yang sama dapat memberi
keputusan berbeda bila ditafsirkan dengan jenis bilangan yang salah.

Pola \texttt{if} memakai satu lompatan bersyarat untuk melewati blok
bila syarat tidak terpenuhi. Pola \texttt{if--else} memakai satu
lompatan bersyarat menuju cabang kedua dan satu \texttt{JMP} untuk
melewati cabang tersebut setelah selesai. \texttt{JMP} pada akhir
cabang pertama mencegah cabang kedua ikut dijalankan.

\begin{rujukan}
Rujukan singkat: Gunawan, Bagian 11.1 dan 11.2; S'to, Bab XIII sampai
sebelum \texttt{LOOP}; praktikum \texttt{bab3.tex} (JMP, CMP, dan
lompatan bersyarat); buku ajar, \texttt{chapters/13} bagian label,
JMP, CMP, dan \texttt{if--else}.
\end{rujukan}

\subsection{Alat dan bahan}

\begin{itemize}
    \item Komputer dengan emu8086 terpasang.
    \item DOSBox dengan TASM, TLINK, dan DEBUG.
    \item Berkas \texttt{code/m04-if.asm} dan
          \texttt{code/m04-if-else.asm}.
    \item Lembar jurnal praktikum.
\end{itemize}

\subsection{Contoh program}

Program pertama memakai satu cabang \texttt{if}. Nilai \texttt{AL}
dibandingkan dengan 5; kalimat dicetak hanya bila \texttt{AL} lebih
besar dari 5 secara tidak bertanda. Bila syarat tidak terpenuhi,
\texttt{JBE} melewati blok pencetakan.

\lstinputlisting[
    caption={Percabangan \texttt{if} dengan \texttt{CMP} dan
    \texttt{JBE}.},
    label={lst:m04-if}
]{code/m04-if.asm}

Program kedua memakai dua cabang. \texttt{JGE} menuju cabang
\texttt{if} bila \texttt{AL} lebih besar atau sama dengan 5 secara
bertanda. Bila tidak, cabang \texttt{else} dijalankan, lalu
\texttt{JMP} melewati cabang \texttt{if}.

\lstinputlisting[
    caption={Percabangan \texttt{if--else} dengan \texttt{JGE} dan
    \texttt{JMP}.},
    label={lst:m04-if-else}
]{code/m04-if-else.asm}

\subsection{Langkah kerja}

\begin{enumerate}
    \item Siapkan folder kerja, lalu salin kedua berkas \texttt{.asm}
          ke dalamnya.
    \item Rakit dan tautkan program pertama. Jalankan kedua perintah
          berikut pada baris terpisah:
          \begin{perintah}
tasm m04-if.asm\\
tlink m04-if.obj
\end{perintah}
          Jalankan hasilnya dan catat kalimat yang tercetak.
    \item Ubah nilai \texttt{AL} pada program pertama dari 7 menjadi
          3, rakit ulang, lalu jalankan. Jelaskan mengapa tidak ada
          kalimat yang tercetak.
    \item Rakit dan tautkan program kedua dengan perintah yang sama,
          lalu jalankan. Catat kalimat yang muncul. Ubah nilai
          \texttt{AL} dari 3 menjadi 7, rakit ulang, dan catat
          perubahannya.
    \item Jalankan \texttt{CMP AL,5} dengan \emph{single step} pada
          emu8086 untuk \texttt{AL=3} dan \texttt{AL=7}. Catat nilai
          \texttt{ZF}, \texttt{CF}, dan \texttt{SF} sesudah
          \texttt{CMP}.
    \item Ubah \texttt{JGE} pada program kedua menjadi \texttt{JA},
          lalu isi \texttt{AL} dengan \texttt{0FFh}. Jalankan dan
          bandingkan hasilnya dengan \texttt{JGE}. Jelaskan mengapa
          kedua lompatan memberi cabang berbeda.
\end{enumerate}

\subsection{Tugas}

\begin{enumerate}
    \item Jelaskan cara kerja \texttt{CMP} dan mengapa kedua
          operandnya tidak berubah.
    \item Jelaskan arti \texttt{ZF}, \texttt{CF}, dan \texttt{SF}
          sesudah sebuah perbandingan.
    \item Diketahui \texttt{AL=0FEh} dan dilakukan
          \texttt{CMP AL,5}. Tentukan apakah \texttt{JA} dan
          \texttt{JG} melompat, lalu jelaskan alasannya.
    \item Jelaskan peran \texttt{JMP} pada pola \texttt{if--else}.
    \item Tuliskan runtunan instruksi untuk pola: bila
          \texttt{AL=5}, cetak kalimat ``sama''; bila tidak, cetak
          kalimat ``tidak sama''. Tentukan lompatan yang dipakai.
\end{enumerate}

\subsection{Hasil yang dikumpulkan}

\begin{enumerate}
    \item Kedua berkas \texttt{.asm} beserta versi modifikasinya.
    \item Tangkapan layar keluaran kedua program pada kedua nilai
          \texttt{AL}.
    \item Tabel nilai \texttt{ZF}, \texttt{CF}, dan \texttt{SF}
          sesudah \texttt{CMP AL,5} untuk \texttt{AL=3},
          \texttt{AL=7}, dan \texttt{AL=0FFh}.
    \item Lembar jurnal berisi jawaban langkah kerja dan tugas.
\end{enumerate}
